\section{Diagnostic Probes and Artifact Catalog}

SID papers now optimize different artifact properties. Learned item indexing
and TIGER-style generative retrieval established the interface~\cite{hua2023indexids,rajput2023tiger}.
Later systems diversify how codes are learned and used: collaborative or
learnable tokenization~\cite{zheng2023lcrec,wang2024letter,liu2024elit},
behavior-semantic or multimodal evidence~\cite{wang2024eager,chen2026syngr},
collision and head-tail objectives~\cite{hu2026quasid,pan2026adasid,yan2026aktrec},
and variable-length, drift, or deployment-facing SID designs~\cite{cheng2026capsid,baikalov2026staleness,ju2026snapchatsid}
all expose different inspection needs.
The literature space is broader than the artifacts that \tool currently
executes. Table~\ref{tab:method-coverage} therefore catalogs the released
resource surfaces rather than repeating the experimental rows. The quantitative
diagnostic evidence appears in Tables~\ref{tab:musical-diagnostic}
and~\ref{tab:mechanism-probes}; Table~\ref{tab:method-coverage} tells a
reviewer what can be inspected in the artifact and where to start.

\begin{table}[t]
\caption{Released resource surfaces.}
\label{tab:method-coverage}
\scriptsize
\setlength{\tabcolsep}{2pt}
\begin{tabularx}{\columnwidth}{@{}>{\raggedright\arraybackslash}p{0.24\columnwidth}>{\raggedright\arraybackslash}X>{\raggedright\arraybackslash}p{0.24\columnwidth}@{}}
\toprule
Resource part & Reviewer-facing check & Entry point \\
\midrule
Adapter contract & New tokenizer exports can be normalized into stable
item-code rows with optional side tables. & README; template \\
Validation gate & Keys, SID levels, joins, and provenance are checked before
diagnostics run. & package verifier \\
D1--D5 probes & Valid artifacts produce utilization, aliasing, alignment,
allocation, and cost reports. & package verifier; CLIs \\
Released adapters & LETTER and LC-Rec item-index files enter the same
contract as local exports. & adapter registry \\
Evidence index & Paper-table numbers are linked to compact snapshots and
commands. &
reproducibility matrix \\
Calibration assets & Control rows and mechanism probes check that D2, D4, and
D5 respond to known stressors. &
evidence snapshots \\
Extension hooks & D6 churn, downstream probes, and D7 trace inputs expose the
extension path. & README; CLIs \\
\bottomrule
\end{tabularx}
\end{table}

\paragraph{Notation.}
Let \(\mathcal{I}\) be the inspected item set and let a tokenizer export a code
sequence \(z(i)=(z_1(i),\ldots,z_L(i))\) for each item \(i\). We write the
level-\(\ell\) prefix as \(p_\ell(i)=(z_1(i),\ldots,z_\ell(i))\) and the full
alias set for code \(c\) as \(A(c)=\{i\in\mathcal{I}:z(i)=c\}\). The probes
below operate on \(z\), with metadata and interactions added only when a
diagnostic needs semantic or collaborative slices.

\paragraph{D1--D5 mapping-level probes.}
D1, \emph{utilization}, reports per-level usage, prefix counts, imbalance
summaries, and dead or underused code indicators. D2, \emph{aliasing},
measures the fraction of items whose full \sid belongs to a non-singleton code
bucket, plus the same quantity for prefixes. This is an item-in-alias-set
profile rather than the number of duplicate code values; causal harm requires
intervention or downstream exposure checks beyond this mapping profile. A
bounded interaction-qualified mechanism probe is included because collision-aware work argues that duplicate-code
assignments are not uniformly harmful, and recent tokenizer-comparison work
shows that SID-level evaluation can diverge from item-level recommendation under
full-code collisions~\cite{hu2026quasid,zhang2026sidreliability}. D3,
\emph{neighborhood alignment}, asks whether prefix neighborhoods recover
train-only item co-occurrence neighbors. For each user, \tool forms item pairs
among bounded train interactions, ranks neighbors by co-occurrence count, keeps
top-\(k\) directed neighbors per item, and reports the fraction whose neighbor
shares \(p_\ell(i)\); the weighted score averages edges, while the mean score
averages items. Metadata purity is auxiliary context, not collaborative
quality. D4, \emph{popularity allocation}, splits items by popularity and measures whether
full-code capacity and prefix structure are allocated differently across head,
mid, and tail items. D5, \emph{structural cost}, reports \sid length, prefix
fan-out, duplicate codes, and trie-like expansion pressure. Measured serving
latency depends on the generator and decoding stack; long, variable-length, or asymmetric SID methods make that distinction
necessary~\cite{xia2026acerec,cheng2026capsid,huang2026asymrec,zhang2026hgrec,wang2026sa2crq,hou2026expressiveness}.

\paragraph{Extension probes.}
D6, \emph{temporal churn}, computes mapping changes between tokenizer
refreshes for drift-aware settings~\cite{feng2026dact,baikalov2026staleness}.
It is implemented as an optional refresh-pair extension.
D7, \emph{generation traces}, covers invalid paths, next-token entropy, and
duplicate generated candidates. D7 requires generator-output tables or beam
traces; the current evidence focuses on the mapping layer.
